\chapter{Development Plan}
\ChapterLabel{DevelopmentPlan}
\index{development plan}
\index{risk}
\index{TBD}

\section{Chapter Responsibility}
\label{Section:DevelopmentPlanResponsibility}
Define ownership, development gates, change control, risks, and the decisions
needed to progress from Week 1 context to an evidenced ADC demonstration.

\noindent\textbf{Expected result:} An actionable proposed plan with named
decision owners and explicit exit evidence.

\section{Governance and Roles}
Jack Marino is the project owner, student author, and proposed hardware,
firmware, integration, and test owner. He resolves technical TBDs and retains
evidence. The course reviewer confirms the academic scope and submission
acceptance; review has not yet occurred. Additional contributors, if any,
will be recorded before ownership changes.

AI assistance is limited to drafting, analysis, and checks under
\ChapterRef{AIAgentGuide} and \ChapterRef{AIImplementationWorkflow}.
One agent acts as aiaCoordinator and aiaDeveloper for this change and performs
a self-review. It cannot approve requirements, safety decisions, or release.
Human review of the resulting pull request remains pending.

\section{Domain Profile Selection}
\label{Section:DevelopmentPlanProfiles}
\small
\begin{tabularx}{\linewidth}{L{3.1cm}L{2cm}Y}
\caption{Project Domain Profile Decisions\label{Table:ProjectDomainProfiles}}\\
\toprule
Profile & Status & Rationale, owner, and date\\
\midrule
\endfirsthead
\toprule
Profile & Status & Rationale, owner, and date\\
\midrule
\endhead
\bottomrule
\endlastfoot
\hyperref[Chapter:EmbeddedSystemsProfile]{Embedded Systems} &
Applicable &
Jack Marino confirmed an embedded system producing ADC readings, with
hardware TBD, on November 10, 2026. The profile adds hardware review,
electrical-interface and timing decisions, target evidence, and power/fault
reviews at the gates below.\\
\end{tabularx}
\normalsize
This is the only supplied domain profile. Any additional safety, security,
or regulatory tailoring depends on the intended environment and the decision
in TBD-05; no certification applicability is asserted.

\section{Milestones and Decision Gates}
Calendar deadlines beyond this submission have not been supplied. The plan
uses ordered gates; Jack Marino agrees dates with the course reviewer under
TBD-06. Every gate below is planned, including Week 1 review.

\small
\begin{tabularx}{\linewidth}{L{2.4cm}Y L{3.1cm}}
\caption{Proposed Development Milestones\label{Table:WeekOneMilestones}}\\
\toprule
Gate & Work and exit evidence & Accountable owner\\
\midrule
\endfirsthead
\toprule
Gate & Work and exit evidence & Accountable owner\\
\midrule
\endhead
\bottomrule
\endlastfoot
G1: Week 1 review &
Introduction, System Context, and Development Plan in a complete document
build; review of scope, TBD ownership, and initial boundaries. &
Jack Marino prepares; course reviewer confirms scope.\\
G2: Requirements and hardware selection &
Resolve TBD-01 through TBD-05 to the extent required for baseline; select
hardware and record documentation, measurable conversion/interface
requirements, and acceptance criteria. &
Jack Marino proposes; course reviewer confirms course scope.\\
G3: Design and bring-up &
Review hardware--software allocation, acquisition timing, output protocol,
resource use, power and fault behavior. Record successful programming and
basic board bring-up evidence. &
Jack Marino.\\
G4: Integration and verification &
Implement the acquisition path, run requirement-mapped tests using a known
input and independent reference, and retain versions, setup, logs, readings,
failures, and dispositions. &
Jack Marino.\\
G5: Demonstration and release review &
Demonstrate the agreed intended use, obtain reviewer feedback, disclose open
defects, and identify the delivered hardware/firmware/document baseline. &
Jack Marino prepares; course reviewer accepts the submission.\\
\end{tabularx}
\normalsize

\section{Week 1 Decision Register}
\label{Section:WeekOneDecisions}
All entries below are open. Jack Marino owns the technical proposals;
the course reviewer confirms scope and academic acceptance where indicated.
The IDs identify questions, not approved requirements.

\small
\begin{tabularx}{\linewidth}{L{1.8cm}Y L{2.7cm}L{2cm}}
\caption{Week 1 TBDs and Decision Owners\label{Table:WeekOneTBDs}}\\
\toprule
ID & Decision needed and impact & Owner & Due gate\\
\midrule
\endfirsthead
\toprule
ID & Decision needed and impact & Owner & Due gate\\
\midrule
\endhead
\bottomrule
\endlastfoot
TBD-01 & Board/processor, integrated or external ADC, board revision,
toolchain, programming method, supply, and hardware availability.
Blocks platform-specific design and bring-up. &
Jack Marino & Before G2 closes\\
TBD-02 & Input source, quantity measured, channel count, voltage range,
reference, grounding, source impedance, bandwidth, conditioning, and
protection. Blocks electrical interface requirements and wiring. &
Jack Marino & Before G2; before connecting hardware\\
TBD-03 & Resolution, sample rate and scheduling, accuracy/tolerance,
latency, acceptable noise, and missed-sample behavior. Blocks measurable
conversion requirements and test acceptance criteria. &
Jack Marino & Before G2 closes\\
TBD-04 & Output interface, receiving/observation tool, numeric format,
units/scaling, framing, and behavior when output cannot keep up.
Blocks the digital interface contract and demonstration method. &
Jack Marino & Before G2 closes\\
TBD-05 & Intended operating environment, startup/reset and fault responses,
power/resource budgets, safety and security concerns, update/service needs,
and any standards or certification relevance. Blocks unsupported safety
claims and final applicability dispositions. &
Jack Marino & At G2; before dependent work\\
TBD-06 & Confirm bench demonstration scope and proposed exclusions,
course deadlines, available equipment, demonstration participants, and
evidence sufficient for academic acceptance. &
Jack Marino proposes; course reviewer confirms & G1 scope; dates before scheduling\\
\end{tabularx}
\normalsize

\section{Risks and Mitigations}
\small
\begin{tabularx}{\linewidth}{L{2.8cm}Y Y}
\caption{Initial Project Risks\label{Table:WeekOneRisks}}\\
\toprule
Risk & Effect & Mitigation and owner\\
\midrule
\endfirsthead
\toprule
Risk & Effect & Mitigation and owner\\
\midrule
\endhead
\bottomrule
\endlastfoot
Hardware or tools unavailable & Bring-up and measurement cannot start. &
Jack Marino checks availability before committing the G2 platform selection;
agrees an alternative with the reviewer if blocked.\\
Input/reference mismatch & Incorrect readings or damaged hardware. &
Jack Marino resolves TBD-02 and checks selected documentation and wiring
before applying signals.\\
Performance unspecified & No objective basis for a demonstration pass. &
Jack Marino resolves TBD-03 and obtains reviewer agreement on acceptance
criteria before implementing performance-dependent behavior.\\
Output throughput mismatch & Missing or stale readings. &
Jack Marino budgets acquisition and output throughput and tests overload
behavior after the interface is selected.\\
Incomplete document build & Submission cannot be reproduced. &
Jack Marino retains all included sources and checks the full build,
references, bibliography, glossary, and index at each document baseline.\\
Template examples mistaken for project facts & False requirements or evidence. &
Jack Marino keeps the Week 1 status explicit and replaces later examples
through reviewed, traceable changes in subsequent course work.\\
\end{tabularx}
\normalsize

\section{Change Control and Evidence}
Changes use a topic branch and focused commits reviewed through a pull request
against \texttt{main}. A proposed change records its reason, affected chapters,
decision-register impact, checks performed, and open issues. Human approval
is recorded only when it occurs. A requirement or interface change triggers
review of affected design, tests, and traceability.

The complete document build uses
\texttt{make -f makefileMacLinux pdf} on a configured Linux/macOS environment
or \texttt{.\textbackslash makefilePowerShell.ps1 -Action pdf} on configured Windows.
\ChapterRef{DocumentGuide} provides the toolchain and manual command sequence.
The repository build workflow retains the generated PDF and logs; its result
is evidence for the associated commit only. Later target testing additionally
records the hardware revision, firmware commit, configuration, equipment,
inputs, expected results, actual results, and defects.

\section{Reviewer Decisions Requested}
At G1, confirm the proposed bench demonstration scope, system boundary,
Embedded Systems profile selection, TBD owners, and staged plan.
Confirm that later template chapters are deferred course work and identify any
additional Week 1 acceptance expectations. At G2, confirm the proposed
requirements and demonstration criteria after the owner supplies hardware and
performance decisions. These decisions remain pending in this submission.
